%% kapa: exam=1 round=26 year=2015 period=2 form=A subject=appraisal_law subject_ko=부동산관계법규 seq=35 number=35
%% source: Q-Net 감정평가사 기출문제 (cst00309) artlSeq=5138001 file=extracted/1008/5138001/제26회_감정평가사 1차_기출/제26회_감정평가사 1차_기출/2교시_감평_A형.pdf
%% conversion: kordoc 4.24.51 -> tools/gen1.py (auto); review=proofread-sonnet; flags=none
\begin{question}{35}
측량ㆍ수로조사 및 지적에 관한 법령상 지적소관청이 관할 등기관서에 그 등기를 촉탁하여야 하는 경우에 해당하지 않는 것은?
\begin{choices}
  \item 축척변경의 경우
  \item 등록사항의 직권정정의 경우
  \item 지번변경의 필요에 의해 지번을 변경한 경우
  \item 행정구역의 개편으로 지번을 새로 부여하는 경우
  \item 신규등록을 이유로 토지이동의 내용을 정리한 경우
\end{choices}
\end{question}
\answer{5}
